\begin{figure}[htbp]
\centering
\begin{tikzpicture}[x=.61cm,y=.61cm]
  \bcuadricula{4}
  \foreach \r/\c in {1/2,2/4,3/1,4/3}{\breina{4}{\r}{\c}}
  \node[font=\small] at (2,-.62) {Complete solution};
  \begin{scope}[shift={(7,0)}]
    \bcuadricula{4}
    \fill[orange,opacity=.24] (0,3) rectangle (4,4);
    \fill[blue,opacity=.15] (1,0) rectangle (2,4);
    \begin{scope}
      \clip (0,0) rectangle (4,4);
      \draw[teal!75!black,line width=.7pt] (-2,7)--(6,-1);
      \draw[red!70!black,dashed,line width=.7pt] (-2,0)--(6,8);
    \end{scope}
    \foreach \r/\c in {2/4,3/1,3/3,4/1,4/3,4/4}{\bmarca{4}{\r}{\c}}
    \breina{4}{1}{2}
    \node[font=\small] at (2,-.62) {One queen revealed};
  \end{scope}
  \begin{scope}[shift={(14,.1)}]
    \foreach \r/\cells in {1/{0,0,1},2/{1,1,0},3/{1,1,1}}{
      \foreach \w [count=\c] in \cells {
        \ifnum\w=1 \fill[teal!16] (\c-1,3-\r) rectangle (\c,4-\r);\fi
        \node[font=\small] at (\c-.5,3.5-\r) {$\w$};
      }
    }
    \draw[step=1,black!45,line width=.3pt] (0,0) grid (3,3);
    \draw[btinta,line width=.6pt] (0,0) rectangle (3,3);
    \foreach \i/\r in {1/2,2/3,3/4}{\node[font=\scriptsize] at (-.3,3.5-\i) {$\r$};}
    \foreach \i/\c in {1/1,2/3,3/4}{\node[font=\scriptsize] at (\i-.5,3.32) {$\c$};}
    \node[font=\small] at (1.5,-.72) {Residual mask};
  \end{scope}
\end{tikzpicture}
\caption{From a solution to the rook problem. After $(1,2)$ is revealed, the dots in the middle are the
ones in the mask on the right, shown with their original labels. Row 2 requires column 4. Columns 1 and 3
can be assigned to rows 3 and 4 in two ways: $\operatorname{per}M=2$. Only the assignment on the left
completes the queen placement. In the other, $(2,4)$ and $(3,3)$ attack each other. Thus $Q(B)=1$.}
\label{fig:mascara-ilustrada}
\end{figure}
